% Auto-generated by generate_narrative_content.py
% Differentiates MiniRocket/MultiRocket from CNN-on-FPGA work
% Insert in Section II (Background) or Related Work

\begin{table}[t]
\centering
\caption{Convolution-Based Classifiers vs.\ Deep CNNs on FPGA}
\label{tab:cnn-vs-conv}
\small
\begin{tabular}{lcc}
\toprule
Property & \textbf{MiniRocket/MultiRocket} & \textbf{CNN (e.g., InceptionTime)} \\
\midrule
Weights & Fixed $\{-1, +2\}$ & Learned (float) \\
Training & Ridge regression only & Backpropagation \\
Kernel count & 84 (combinatorial) & 100s--1000s (learned) \\
Parameters & $<$10K (classifier only) & 100K--10M+ \\
DSP usage & 0 (ternary multiply) & High (float MAC) \\
Accuracy & Competitive & State-of-the-art \\
Feature count & 840--50K (pooling) & Implicit (feature maps) \\
\bottomrule
\end{tabular}
\vspace{-0.5em}
\end{table}

% Accompanying paragraph:
Unlike CNN-based FPGA accelerators~\cite{venieris2018toolflows,guo2017angel}
that focus on mapping learned floating-point convolutions to DSP arrays,
our work targets \emph{convolution-based classifiers} where kernels are
\emph{fixed at design time}.
This distinction is critical: the ternary weight structure ($\{-1, +2\}$)
enables a binary floating-point (bFP) multiplier that uses only XNOR and
addition, consuming \textbf{zero DSP slices}.
The classifier is a simple Ridge regression, not a multi-layer network,
so the entire inference pipeline maps to a single-pass dataflow architecture
without the layer-by-layer scheduling complexity of CNN accelerators.
